Ce chapitre est le socle calculatoire de la partie : valeur absolue, identités remarquables,
puissances, sommes, récurrence. Rien n'y est difficile pris isolément, et presque tout sert
plus loin.

Deux résultats sortent du lot. L'irrationalité de $\sqrt{2}$ ferme la chaîne d'inclusions
$\mathbb{N} \subset \mathbb{Z} \subset \mathbb{D} \subset \mathbb{Q} \subset \mathbb{R}$, dont
chaque maillon est ici montré strict par un contre-exemple explicite. Et le \emph{principe de
récurrence} apparaît pour la première fois comme un énoncé démontré, et non comme une règle
du jeu : en Lean, il est la forme même de la définition des entiers naturels.

Un énoncé du programme n'est pas traité : le développement décimal d'un rationnel est
périodique, et réciproquement. Il demande de manipuler des suites de chiffres et leur
périodicité, ce que le chapitre ne fait pas.
